\documentclass{article}
\usepackage[utf8]{inputenc}
\usepackage[T1]{fontenc}
\usepackage{geometry}
\geometry{
    a4paper,
    left=2.5cm,
    right=2.5cm,
    top=2.5cm,
    bottom=2.5cm
}
\usepackage{graphicx}
\usepackage{float}
\usepackage{hyperref}
\usepackage[english]{babel}
\usepackage{amsmath}
\usepackage{amssymb}
\usepackage{array}
\usepackage{booktabs}
\usepackage{caption}
\usepackage{subcaption}
\usepackage{enumitem}
\usepackage{fancyhdr}
\usepackage{titlesec}
\usepackage{url}
\usepackage{textcomp}
\usepackage{xcolor}
\usepackage{listings}
\usepackage{multirow}
\usepackage{longtable}
\usepackage{tabularx}
\usepackage{times}

% Set up page style
\pagestyle{fancy}
\fancyhf{}
\rfoot{\thepage}
\renewcommand{\headrulewidth}{0pt}

% Title formatting
\titleformat{\chapter}[display]{\normalfont\bfseries}{\chaptertitle\thechapter}{0pt}{\Huge}
\titlespacing*{\chapter}{0pt}{50pt}{40pt}

% Document info
\title{Template Document}
\author{Academic Support Agent}
\date{\today}

\begin{document}
\maketitle

\fancyhf{}
\pagestyle{fancy}
\fancyfoot[R]{\footnotesize \thepage}
\lhead{\footnotesize \sc Onderzoekscompetenties wiskunde 6ECWW-6GRWW-6LAWW}
\rhead{\footnotesize \sc Schooljaar 2017-2018}

%Start Hoofding%%%%%%%%%%%%%%%%%%%%%%%%%%%%%%%%%%%%%%%%%%%%%%%%%%%%%%%%%%%%%%%%%%%%%%%%%%%%

\scalebox{0.96}[1]{\begin{tabular}[t]{||c||ll||}
\hline  \hline
\raisebox{-0.5cm}{\includegraphics[scale=0.4]{logoschool.eps}} &
\begin{tabular}{ll|ll}
Ond:&... & Groepsleden:& ........ \\
& & & ........ \\
& & & .............. \\
& & & .............. \\
& & &\\
\end{tabular}
& \\
\sc Hoogstraat 30 & \raisebox{-0.3cm}{\kader{yellow}{\textbf{Verslag onderzoekscompetenties wiskunde (Onderwerp .....)}}} &  \\
\sc 9700 Oudenaarde&   & \\
\hline \hline
\end{tabular}\\[1ex]}
%Einde Hoofding%%%%%%%%%%%%%%%%%%%%%%%%%%%%%%%%%%%%%%%%%%%%%%%%%%%%%%%%%%%%%%%%%%%%%%%%%%%%



\newpage
\section{test}
Hier kan je beginnen schrijven...\cite{test}%deze laatste regel is een link naar je bronvermelding

\bigskip
\begin{wrapfigure}{r}{0.25\textwidth}
    \centering
  	\includegraphics[scale=0.3]{zonderGRM.JPG}
\end{wrapfigure}

Zo kan je dus figuren toevoegen in je bestand en je tekst eromheen ``wrappen''. Via de knop ``PROJECT'' bovenaan zie je welke files er in dit document staan.

\subsection{test}

En zo typ je wiskunde:
\begin{itemize}
\item Stelsels:
$$\left\{\begin{array}{rcl}
f(x)&=& 6\sin\left[\dfrac{\pi}{300}(x-323)\right]+36\\
g(x)&=& \dfrac{36}{323}x
\end{array}\right.$$

\item Matrices:
$\bordermatrix{\text{ }&K_1&K_2&K_3&K_4\cr
                L_a& 1 &  0  & 0 & 1\cr
                L_b& 0  &  1 & 1 & 0\cr
                L_c& 1 & 0 & 1 & 0\cr
                L_d& 0 & 1 & 0 & 1}$ of
$\begin{pmatrix}
1 & 0 & 0 & | & 1 & 0\\
0 & 1 & 0 & | & 1 & 0\\
1 & 0 & 1 & | & 0 & 0\\
- & - & - & | & - & -\\
1 & 0 & 0 & | & 0 & 1\\
0 & 1 & 0 & | & 0 & 1
\end{pmatrix}$

\item Besluiten in fancy kadertjes:

\kaderm{yellow}{17cm}{De sinusfunctie $g(x)$ met bovenstaande waarden heeft als voorschrift:\begin{center}$g(x)= 7\sin\left[\dfrac{\pi}{19}(x-9)\right]+7$\end{center}}
\end{itemize}

\begin{thebibliography}{1}
% hier komen je bronnen. 
\bibitem{test}
\url{http://www.google.be}



\end{thebibliography}
\end{document}